Wir verwenden
[[Differenzierbare Funktion/Graph/Hyperfläche/Weingartenabbildung/Fakt|Fakt]],
wobei im Nullpunkt der Vorfaktor gleich $1$ ist und somit die
\definitionsverweis {Weingartenabbildung}{}{}
direkt durch die
\definitionsverweis {Hesse-Matrix}{}{}
beschrieben wird. Die Hesse-Matrix der Funktion ist

\mathdisp {\begin{pmatrix}  2a  &   c   \\ c &  2b \end{pmatrix}} { , }

das
\definitionsverweis {charakteristische Polynom}{}{}
davon ist

\mavergleichskettealign
{\vergleichskettealign
{ (t-2a)(t-2b)-c^2
}
{ =} { t^2 -2(a+b)t +4ab -c^2
}
{ =} { \left(  t-(a+b)  \right)^2 - (a+b)^2 +4ab -c^2
}
{ =} { \left(  t-(a+b)  \right)^2  - a^2 -b^2 +2ab -c^2
}
{ =} { \left(  t-(a+b)  \right)^2  - (a-b)^2 -c^2
}
}
{}
{}{.}
Die Eigenwerte
\zusatzklammer {und dies sind die Hauptkrümmungen} {} {}
sind also

\mavergleichskettedisp
{\vergleichskette
{ \lambda_{1,2}
}
{ =} { \pm \sqrt{ (a-b)^2+c^2 } +a+b
}
{ } {
}
{ } {
}
{ } {
}
}
{}{}{.}
Bei

\mavergleichskette
{\vergleichskette
{ a
}
{ = }{ b
}
{  }{
}
{    }{
}
{   }{
}
}
{}{}{}
und

\mavergleichskette
{\vergleichskette
{ c
}
{ = }{ 0
}
{  }{
}
{    }{
}
{   }{
}
}
{}{}{}
ist die Weingartenabbildung Multiplikation mit $2a$ und jeder Vektor ist ein Eigenvektor, sei dies jetzt ausgeschlossen. Für

\mavergleichskettedisp
{\vergleichskette
{ \lambda_{1}
}
{ =} { \sqrt{ (a-b)^2+c^2 } +a+b
}
{ } {
}
{ } {
}
{ } {
}
}
{}{}{}
ist der Kern der Matrix

\mathdisp {\begin{pmatrix}  \sqrt{(a-b)^2 +c^2} + b-a  &   -c  \\  -c &  \sqrt{(a-b)^2 +c^2} + a-b \end{pmatrix}} {  }

zu bestimmen, dieser ist

\mathdisp {\R \begin{pmatrix}  \sqrt{(a-b)^2 +c^2} + a-b  \\c   \end{pmatrix}} { . }

Eine Hauptkrümmungsrichtung ist also
\mathl{\begin{pmatrix}  \sqrt{(a-b)^2 +c^2} + a-b  \\c   \end{pmatrix}}{} mit der Hauptkrümmung
\mathl{\sqrt{ (a-b)^2+c^2 } +a+b}{.}

Entsprechend ergibt sich die Hauptkrümmungsrichtung
\mathl{\begin{pmatrix}  -c \\ \sqrt{(a-b)^2 +c^2} + a-b    \end{pmatrix}}{} mit der Hauptkrümmung
\mathl{-\sqrt{ (a-b)^2+c^2 } +a+b}{.}

 [[Kategorie:Latexseite]]
